\section{Output Targets}
\label{sec:outputs}

The rendering pipeline produces a deterministic \emph{Scene}---a backend-agnostic
description of every drawing primitive, visual treatment, and effect request---from an IR
document and theme. The Scene is the stable contract between the layout engine and any
output format. Rendering backends consume the Scene and emit their respective outputs.
This architecture replaces the prior design in which SVG was the primary format from
which all other formats were derived.

\begin{figure}[h]
\centering
\begin{Verbatim}[frame=single,fontsize=\small]
  IR Document + Theme
         |
         v
  Layout Engine (sec.5 -- deterministic six-phase pipeline)
         |
         v
    Scene / Render IR  <-- byte-deterministic; backend-agnostic
   /         |          \
  v          v           v
SVG       Raster       PPTX
Backend   Backend      Native-Shape
(sec.7.3) (sec.7.4)    Backend (sec.7.5)
 |  \       / \             |
SVG PDF   PNG  PDF(art)   PPTX
          \                 |
           HTML/Interactive (sec.7.6) -- wraps SVG or raster
\end{Verbatim}
\caption{Revised output pipeline: Scene/Render IR as root; SVG, Raster, and PPTX as
  pluggable backends; PNG and PDF are produced per backend, not universally from SVG}
\label{fig:output-pipeline}
\end{figure}

\paragraph{Why SVG is one beloved backend, not the universal root.}
SVG is not incapable of visual effects. \texttt{feGaussianBlur}, \texttt{feDropShadow},
\texttt{feTurbulence}, and \texttt{feColorMatrix} filter primitives exist in SVG and can
approximate glow, drop shadows, and noise textures. The architectural problem is different.
First, SVG filter rendering is \emph{not deterministic or consistent} across SVG renderers
and browsers: identical filter markup produces visually different results in Inkscape,
Chrome, Safari, Firefox, and resvg. This breaks the project's byte-determinism principle
for any filter-bearing output. Second, effects expressed as SVG filters do not survive
faithfully when rasterising to PNG or converting to PDF, and they have no native equivalent
in PPTX's shape model. Third---and most consequentially---anchoring every output target to
SVG's semantics means the most expressive targets are constrained by the least expressive:
clean, deterministic vector geometry is SVG's strength, but it is not art effects' domain.
The Scene as root resolves this: it describes \emph{what} to draw, including effect
requests with explicit fallback policies, independently of \emph{how} any backend draws it.

\subsection{Scene / Render IR}
\label{sec:scene-ir}

The Scene is the immutable, byte-deterministic record produced by the six-phase layout
pipeline (Section~\ref{sec:layout-algorithm}). It is the single handoff point between the
layout engine and every rendering backend.

\subsubsection{Scene Model}
\label{sec:scene-model}

The Scene consists of three top-level structures: a \textbf{canvas descriptor}, an
\textbf{ordered element list}, and an \textbf{effect registry}. The following is a
data-level description, not code:

\begin{Verbatim}[frame=single,fontsize=\small]
Scene {
  canvas:   { width, height, background_color }
  elements: [ DrawingPrimitive ]   -- ordered; painting order preserved
  effects:  { effect_id -> EffectDefinition }
  meta:     { theme_id, fidelity_tier, scene_hash(SHA-256) }
}

DrawingPrimitive := one of:
  Rect    { id, x, y, width, height, fill, stroke, stroke_width,
            corner_radius, opacity, clip_id?, effect_refs[] }
  Polygon { id, vertices[(x,y)...], fill, stroke, stroke_width, opacity }
  Line    { id, x1, y1, x2, y2, stroke, stroke_width, stroke_style,
            opacity }
  Text    { id, x, y, content, font_family, font_size, font_weight,
            color, anchor, clip_id? }
  Path    { id, d_commands[(op,args)...], fill, stroke, stroke_width,
            opacity }
  Image   { id, x, y, width, height, data_uri }  -- embedded raster
  Group   { id, children[DrawingPrimitive], clip_rect? }

EffectDefinition := one of:
  Glow         { color, radius_px, intensity, fallback_policy }
  DropShadow   { offset_x, offset_y, blur_radius, color, opacity,
                 fallback_policy }
  GradientFade { direction, color_stops[(offset,color)...], fade_px }
  NoiseTexture { seed, scale, turbulence_type, opacity, color,
                 fallback_policy }
  CloudLayer   { seed, scale, octaves, opacity, color_stops,
                 fallback_policy }
  Bloom        { radius_px, threshold, intensity, fallback_policy }
\end{Verbatim}

\paragraph{Effect requests and fallback policies.}
Each EffectDefinition carries a \texttt{fallback\_policy} that instructs a backend what to
do when it cannot render the effect natively:

\begin{center}
\small
\begin{tabular}{ll}
\toprule
\textbf{Policy value} & \textbf{Meaning} \\
\midrule
\texttt{approximate}  & Use the closest available native mechanism \\
\texttt{omit}         & Silently omit the effect; geometry unchanged \\
\texttt{embed-raster} & Rasterise the affected layer; embed as \texttt{Image} \\
\texttt{error}        & Abort with a descriptive diagnostic message \\
\bottomrule
\end{tabular}
\end{center}

Themes declare effects and their fallback policies
(Section~\ref{sec:schema-fidelity}~and~\ref{sec:fidelity-tiers}). The layout engine does
not choose backends; it records effect requests faithfully in the Scene. The backend
resolves fallback when it cannot fulfil a request.

\subsubsection{Scene Determinism}
\label{sec:scene-determinism}

The Scene is byte-deterministic: two executions of the layout pipeline on identical IR and
theme produce bitwise-identical Scene records. The \texttt{scene\_hash} (SHA-256 of the
serialised element list and effect registry) is included in Scene metadata and is
reproducible across runs.

\paragraph{Layered determinism contract.}
The determinism guarantee applies at three layers, each with different strength:

\begin{enumerate}
  \item \textbf{Scene geometry---always byte-deterministic.} All coordinates, dimensions,
        label positions, and element ordering are bitwise-reproducible. This is guaranteed
        unconditionally by the pure pipeline contract in
        Section~\ref{sec:rendering-determinism}.

  \item \textbf{Per-backend output---deterministic given a pinned backend version.}
        Each backend is a deterministic function of (Scene, backend-version,
        effect-seeds). Effect seeds for all procedural effects (\texttt{NoiseTexture},
        \texttt{CloudLayer}) are derived from the \texttt{scene\_hash} concatenated with
        the \texttt{effect\_id}; no random state is consumed. The backend version must
        be recorded in the output's metadata.

  \item \textbf{Cross-backend pixel identity---explicitly not promised.}
        The SVG backend and the Raster backend, given the same Scene, are permitted and
        expected to produce visually different outputs. This is correct: the raster backend
        renders art effects that the SVG backend cannot. Two outputs from the same Scene
        but different backends are both correct renderings at different fidelity tiers.
        Cross-backend tests use per-backend golden images~\cite{golden-image-testing},
        not cross-backend pixel equality.
\end{enumerate}

\subsection{SVG Backend}
\label{sec:output-svg}

\textbf{Role:} The clean, scalable, print and web default. SVG is the recommended first
implementation target and the most widely compatible vector format. It is one beloved
backend with a well-defined capability ceiling---not the universal root.

\paragraph{Benefits.}
\begin{itemize}
  \item \textbf{Deterministic at Tier~0 and Tier~1.} Scene geometry and Tier-0/1 effects
        (solid fills, linear and radial gradients, patterns, opacity) serialise to SVG
        with full byte-determinism.
  \item \textbf{Scalable.} Renders crisply at any resolution; A0 conference poster to
        screen thumbnail.
  \item \textbf{Inspectable.} XML text format; human-readable, version-controllable with
        meaningful diffs, and inspectable in browser developer tools.
  \item \textbf{Embeddable.} Inline HTML, Markdown image references (\texttt{![]()}),
        \LaTeX\ \texttt{\textbackslash{}includegraphics}, or direct browser rendering.
  \item \textbf{Accessible.} Live \texttt{<text>} nodes support copy-paste, search, and
        screen readers.
\end{itemize}

\paragraph{Limitations and honest scope.}
\begin{itemize}
  \item SVG filters (\texttt{feGaussianBlur}, \texttt{feDropShadow},
        \texttt{feTurbulence}) exist and can approximate glow, shadows, and noise.
        However, filter rendering is \emph{not deterministic across SVG renderers}: the
        same filter markup produces visually different results in Inkscape, Chrome, Safari,
        Firefox, and resvg. For this reason, SVG filters are restricted to a ``safe
        filter'' set (Table~\ref{tab:svg-safe-filters}) at Tier~2 with a determinism
        caveat; all Tier-3 effects require the Raster backend.
  \item Text rendering varies across SVG viewers when system fonts are used. Mitigation:
        embed font files or convert text to paths for archival exports.
  \item Art effects (glow/bloom, cloud textures, volumetric atmospheres, complex noise)
        are not achievable at Tier-3 fidelity in SVG. A Tier-3 theme targeting the SVG
        backend triggers effect fallback policies
        (Section~\ref{sec:scene-ir}).
  \item Complex diagrams with many activities produce proportionally large SVG files
        (no rasterisation compression).
\end{itemize}

\begin{table}[h]
\centering
\small
\begin{tabular}{lll}
\toprule
\textbf{Effect} & \textbf{SVG construct} & \textbf{Determinism} \\
\midrule
Linear gradient  & \texttt{<linearGradient>} in \texttt{<defs>} & Fully deterministic \\
Radial gradient  & \texttt{<radialGradient>} in \texttt{<defs>} & Fully deterministic \\
Pattern (hatch)  & \texttt{<pattern>} in \texttt{<defs>}         & Fully deterministic \\
Opacity          & \texttt{opacity} attribute                     & Fully deterministic \\
Drop shadow      & \texttt{feDropShadow} composite                & Non-deterministic (caveat) \\
Soft glow        & \texttt{feGaussianBlur + feComposite}          & Non-deterministic (caveat) \\
Noise / texture  & \texttt{feTurbulence}                          & Non-deterministic (caveat) \\
\bottomrule
\end{tabular}
\caption{SVG effect constructs and determinism status; ``Non-deterministic (caveat)'' means
  the construct exists but renders inconsistently across SVG renderers and browsers}
\label{tab:svg-safe-filters}
\end{table}

\paragraph{Implementation strategy.}
\begin{enumerate}
  \item Consume the Scene element list in painting order. Serialise each primitive to its
        SVG equivalent:
        \texttt{Rect}$\to$\texttt{<rect>};
        \texttt{Polygon}$\to$\texttt{<polygon>};
        \texttt{Line}$\to$\texttt{<line>};
        \texttt{Text}$\to$\texttt{<text>};
        \texttt{Path}$\to$\texttt{<path>};
        \texttt{Group}$\to$\texttt{<g>};
        \texttt{Image}$\to$\texttt{<image>}.
  \item Apply SVG serialisation invariants:
        \begin{itemize}
          \item All coordinates rounded to 2 decimal places (round-half-up).
          \item All colour values as 6-digit hex (\texttt{\#RRGGBB}); no named colours,
                no \texttt{rgb()} syntax.
          \item All dimension values as integers for rectangles and polygons; 2-decimal
                for text positions.
          \item Font references via \texttt{<style>} or \texttt{@font-face} in
                \texttt{<defs>}; no inline \texttt{style=} font declarations.
          \item Pattern definitions in \texttt{<defs>}, referenced by stable ID
                (\texttt{url(\#tc-pattern-diagonal-hatch)}).
          \item Diamond vertices computed directly; no
                \texttt{transform="rotate(...)"} on individual elements.
          \item Element IDs use prefix \texttt{tc-} followed by entity \texttt{id}.
          \item Elements serialised in painting order: sections, gridlines, activity bars,
                milestones, annotations, labels.
          \item Two SVG outputs from identical Scenes must be byte-identical.
        \end{itemize}
  \item For Tier-2 effects: emit the safe-filter SVG constructs from
        Table~\ref{tab:svg-safe-filters}; record the backend version in a
        \texttt{<desc>} element and include a
        \texttt{<!-- tc:determinism-caveat:filter -->} comment in the SVG header.
  \item For Tier-3 effects: apply the fallback policy from the Scene. If
        \texttt{omit}, skip the effect entirely. If \texttt{embed-raster}, delegate the
        affected layer group to the Raster backend and embed the result as
        \texttt{<image>}.
\end{enumerate}

\paragraph{Text-as-paths option.}
An optional \texttt{--text-as-paths} flag converts all \texttt{<text>} elements to
\texttt{<path>} elements using pre-computed glyph outlines from embedded font files.
This eliminates font-dependent rendering differences at the cost of SVG searchability and
file size. It is not the default.

\paragraph{PDF via SVG backend.}
For vector/consulting targets (Tier~0 and Tier~1), PDF is produced by converting the SVG
backend output using a programmatic library. Recommended: \texttt{svg2pdf} (Rust,
Apache~2.0) or \texttt{cairosvg} (Python, MIT). \textbf{Not} LibreOffice, headless Chrome,
or Puppeteer (too heavyweight; introduces rendering non-determinism). Fonts are embedded
as Type~1/CFF; PDF metadata is set deterministically from IR fields
(\texttt{/Title}$\leftarrow$\texttt{metadata.title};
\texttt{/Author}$\leftarrow$\texttt{metadata.author};
\texttt{/CreationDate}$\leftarrow$\texttt{metadata.created}, not the system timestamp).
The PDF document ID is a SHA-256 hash of the SVG content bytes, truncated to 16 bytes.

\subsection{Raster Backend --- High-Fidelity Art Effects}
\label{sec:output-raster}

\textbf{Role:} Full art effects---glow, bloom, soft shadows, volumetric cloud and
atmosphere layers, noise textures, gradient meshes, blur, and per-pixel compositing.
This is the differentiator for presentation showcase and keynote use cases.

\textbf{Candidate implementations:} Skia~\cite{skia} (C++ with Rust bindings; used in
Chrome, Flutter, and Android; permissive licence), headless HTML
Canvas~\cite{webgl} (Node.js at a fixed, pinned version), or a WebGL-based headless
pipeline~\cite{webgl} for GPU-accelerated art effects.

\paragraph{Benefits.}
\begin{itemize}
  \item Renders every Scene effect class natively: glow, bloom, drop shadows, soft
        shadows, cloud and atmosphere textures via noise/turbulence compositing, gradient
        meshes, per-pixel alpha compositing, and blur at any kernel radius.
  \item Output is a raster bitmap (PNG primary; also packaged as PDF for art-target print).
  \item Well-understood determinism control surfaces: seed pinning, fixed sampling
        parameters, pinned library version.
\end{itemize}

\paragraph{Determinism strategy for art effects.}
Art effects are inherently parameterised and continuous. Determinism is preserved by:
\begin{enumerate}
  \item \textbf{Pinned renderer version.} The raster backend records its version
        (\texttt{skia-v87}, \texttt{canvas-20.x}, etc.) in the output metadata. A
        given version produces identical output for identical inputs; the version is an
        explicit parameter of the determinism contract.
  \item \textbf{Fixed effect seeds.} All procedural effects (\texttt{NoiseTexture},
        \texttt{CloudLayer}) derive seeds from \texttt{scene\_hash} concatenated with the
        \texttt{effect\_id}. No random state is consumed; seeds are deterministic functions
        of the Scene.
  \item \textbf{Fixed sampling and rounding.} Anti-aliasing uses MSAA at $4\times$ (fixed);
        sub-pixel AA is disabled (produces different results across LCD orientations).
        Font rendering uses pre-computed glyph outlines. PNG compression: deflate level~6
        (not level~9: marginal savings, different byte sequences across zlib versions).
  \item \textbf{Golden-image testing.} Reference PNGs are generated per backend version
        and stored as test fixtures~\cite{golden-image-testing}. CI compares new output
        to the golden image using a strict pixel-equality assertion (zero tolerance for
        geometry; configurable tolerance for anti-aliasing edges, default 0 pixels).
\end{enumerate}

\paragraph{Limitations.}
\begin{itemize}
  \item Raster output is fixed resolution; upscaling produces pixelation. Default DPI:
        150 for screen/web; 300 for print-quality.
  \item The backend introduces a runtime dependency (Skia or headless browser) that
        increases distribution size. Recommended: ship as an optional plugin activated
        for Tier-2 and Tier-3 themes only (Section~\ref{sec:distribution}).
  \item Cross-backend pixel identity with the SVG backend is explicitly \emph{not}
        promised; see the layered determinism contract (\S\ref{sec:scene-determinism}).
\end{itemize}

\paragraph{PDF via Raster backend.}
For art-target PDF (Tier-3 themes), PDF is produced by embedding the raster output as a
full-page image in a PDF/A envelope. Vector PDF (SVG backend) and art PDF (Raster backend)
are distinct outputs selectable via a CLI flag (\texttt{--pdf-backend=vector|raster}).

\subsection{PPTX --- Native-Shape Backend}
\label{sec:output-pptx}

\textbf{Format:} Office Open XML Presentation (ISO/IEC~29500).
\textbf{Library:} \texttt{python-pptx} (MIT licence)~\cite{python-pptx}.

\paragraph{Architecture: native shapes with embedded raster art layers.}

The PPTX backend maps Scene primitives to native PowerPoint shapes wherever possible,
enabling think-cell\,\cite{thinkcell}-comparable editability, and falls back to embedded
raster layers only for effects that have no native PowerPoint equivalent.

\begin{table}[h]
\centering
\small
\begin{tabularx}{\textwidth}{lXXX}
\toprule
\textbf{Strategy} & \textbf{Fidelity} & \textbf{Editability} & \textbf{Implementation} \\
\midrule
Image embedding &
  Perfect fidelity to raster output &
  None: opaque image element &
  Trivial (one API call per slide) \\[6pt]
Native shapes &
  Geometry faithful; effects approximated or omitted &
  Full: bars, diamonds, and labels are movable, recolourable MSO shapes &
  High: EMU conversion, MSO shape API, hatch XML \\[6pt]
Hybrid (recommended) &
  High for Tier-0/1 geometry; approximated for Tier-2 effects; rasterised for Tier-3 art layers &
  Partial: native shapes editable beneath opaque art-layer overlays &
  Medium \\
\bottomrule
\end{tabularx}
\caption{PPTX output strategies and trade-offs}
\label{tab:pptx-strategies}
\end{table}

\paragraph{Native PowerPoint effects: exploiting what exists.}
PowerPoint's shape model includes native effect properties that satisfy Tier-2 effect
requests without embedded images. The PPTX backend exploits them as follows~\cite{ooxml}:

\begin{itemize}
  \item Scene \texttt{Glow} $\to$ PowerPoint shape \texttt{Glow} property
        (\texttt{a:glow} element), matching colour and scaled radius. Native; editable.
  \item Scene \texttt{DropShadow} $\to$ PowerPoint \texttt{OuterShadow}
        (\texttt{a:outerShdw}), matching offset, blur, and colour. Native; editable.
  \item Scene \texttt{Bloom} (heavy glow) $\to$ approximated as a larger, lower-opacity
        \texttt{Glow}; fallback policy \texttt{approximate}.
  \item Scene \texttt{CloudLayer} / \texttt{NoiseTexture} $\to$ no native PPTX
        equivalent; rasterised to PNG at 150~DPI and embedded as a transparent image shape
        above the native shapes. Policy: \texttt{embed-raster}. Native shapes beneath
        remain fully editable.
\end{itemize}

\paragraph{Implementation strategy.}
\begin{enumerate}
  \item \textbf{Slide dimensions.} Set to theme canvas width and computed height.
        Coordinate conversion: 1~logical pixel~$= 914.4$~EMU (96~DPI logical);
        round-half-up. This conversion is deterministic.
  \item \textbf{Canvas background.} Full-slide \texttt{RECTANGLE} in
        \texttt{theme.canvas.background\_color}.
  \item \textbf{Track headers.} \texttt{MSO\_AUTO\_SHAPE\_TYPE.RECTANGLE} with
        \texttt{TextFrame}. Font from theme typography. Transparent or light fill per theme.
  \item \textbf{Activity bars.} \texttt{ROUNDED\_RECTANGLE} with corner radius from
        \texttt{theme.activity.bar\_radius}. Fill from resolved status/category style.
        Diagonal-hatch: second semi-transparent overlay with MSO hatch pattern fill.
  \item \textbf{Milestone diamonds.} \texttt{MSO\_AUTO\_SHAPE\_TYPE.DIAMOND} shape.
  \item \textbf{Labels.} \texttt{TextFrame} inside shapes; free-floating \texttt{TextBox}
        for labels placed outside bars.
  \item \textbf{Axis.} Rectangle with tick-label \texttt{TextBox} elements; gridlines as
        thin \texttt{Line} shapes. Today marker as a vertical \texttt{Line} shape.
  \item \textbf{Tier-2 effects.} Apply native PowerPoint effect properties as described
        above. Record the backend version in slide notes metadata.
  \item \textbf{Tier-3 art layers.} Invoke the Raster backend for the affected Scene
        groups; embed results as transparent \texttt{<pic>} shapes above the native layer.
  \item \textbf{Internal IDs.} Derive all \texttt{r:id} values from entity \texttt{id}
        fields (SHA-256 prefix), not auto-generated, for byte-stable output across runs.
\end{enumerate}

\paragraph{Think-cell comparison.}
Think-cell achieves the highest-quality editable output: every element is a native
PowerPoint shape at the XML level, fully fidelity in font, colour, and position~\cite{thinkcell}.
Its approach is PowerPoint-locked and requires a proprietary COM add-in. Our target is to
approach think-cell editability for the Tier-0/Tier-1 element vocabulary (bars, diamonds,
text boxes) while retaining full pipeline determinism and a text-format IR. For Tier-2
themes (Executive, Product), the hybrid strategy produces editable bars and diamonds with
native drop shadows and glow---no embedded rasters needed. For Tier-3 (Showcase), art
layers are embedded rasters atop the fully editable geometric scaffold.

\paragraph{Limitations.}
\begin{itemize}
  \item Diagonal-hatch patterns require the \texttt{pptx.oxml} XML interface, not a
        first-class \texttt{python-pptx} API.
  \item Gradient fades at approximate-date bar edges fall back to the image overlay
        strategy.
  \item PPTX files must be set to a fixed pixel dimension; viewers on different DPI
        settings scale the output proportionally.
\end{itemize}

\subsection{HTML --- Interactive Backend}
\label{sec:output-html}

\textbf{Format:} HTML5 with inline SVG (Tier~0/1 path) or with embedded raster frame
(Tier~2/3 path).

\paragraph{Benefits.}
\begin{itemize}
  \item \textbf{Zero-install preview.} Any browser can display the file. Ideal for the
        VS~Code extension live-preview pane and for sharing timelines via URL.
  \item \textbf{Interactive.} JavaScript hover handlers display \texttt{description}
        and \texttt{metadata} fields from the IR as tooltips, without altering geometry.
  \item \textbf{Accessible.} SVG \texttt{<text>} nodes remain live; labels are
        searchable and selectable. \texttt{<title>} elements on shapes are readable by
        assistive technology.
  \item \textbf{MCP-agent integration.} The MCP server (Section~\ref{sec:agent-integration})
        can return an HTML URL as the rendering result, allowing agents to preview timelines
        in a browser tab without file-system access.
\end{itemize}

\paragraph{Two HTML rendering paths.}
\begin{enumerate}
  \item \textbf{SVG-backed HTML (Tier~0/1, default).} Take the SVG backend output
        as-is. Wrap in an HTML5 document shell: \texttt{<!DOCTYPE html>},
        \texttt{<meta charset="utf-8">}, \texttt{<title>} from \texttt{metadata.title},
        minimal CSS for page background and centre alignment. Inject the SVG inline (not
        as \texttt{<img src="...svg">}) to allow JavaScript access to SVG DOM nodes.
        Attach a lightweight JS event listener (no framework dependency) reading
        \texttt{data-description} and \texttt{data-id} attributes from SVG elements
        and displaying a tooltip \texttt{<div>} on hover. A \texttt{--no-js} flag
        produces geometry-identical HTML with the JavaScript block omitted.
  \item \textbf{Canvas-backed HTML (Tier~2/3, art effects).} Embed a raster-backend PNG
        inside an HTML shell with an overlaid interactive layer. The PNG provides
        art-effect fidelity; the JS layer adds interactivity via a companion JSON map of
        entity bounding boxes (derived from the Scene) to description strings.
\end{enumerate}

\paragraph{Limitations.}
\begin{itemize}
  \item The interactive state (hover, click, tooltip visibility) is not deterministic and
        is not part of the rendering contract. The underlying geometry is identical to the
        chosen backend output.
  \item Browser-to-browser rendering differences in SVG filter effects may produce minor
        visual variations; these are browser environment differences, not Timeline Compiler
        defects.
  \item Not suitable for archival or print without stripping the JavaScript layer.
\end{itemize}

\subsection{Backend Capability and Fidelity-Tier Table}
\label{sec:backend-capability}

The following table summarises which effect classes each backend supports natively,
approximates, or cannot perform. Fallback behaviour applies the policy declared in the
Scene's EffectDefinition.

\begin{table}[h]
\centering
\small
\begin{tabularx}{\textwidth}{lXXX}
\toprule
\textbf{Effect Class} & \textbf{SVG Backend} & \textbf{Raster Backend} & \textbf{PPTX Backend} \\
\midrule
Solid fills            & Native               & Native                  & Native \\
Linear/radial gradient & Native               & Native                  & Native \\
Gradient mesh          & Approximated         & Native                  & Approximated \\
Diagonal hatch         & Native (pattern)     & Native                  & Native (hatch fill) \\
Opacity/transparency   & Native               & Native                  & Native \\
Clip paths             & Native               & Native                  & Native (limited) \\
Drop shadow            & Non-det.\ filter     & Native                  & Native shape effect \\
Soft glow              & Non-det.\ filter     & Native                  & Native shape effect \\
Bloom / light scatter  & Not available        & Native                  & Approximated \\
Diffuse soft shadow    & Non-det.\ filter     & Native                  & Approximated \\
Blur                   & Non-det.\ filter     & Native                  & Not available \\
Noise / turbulence     & Non-det.\ filter     & Native                  & Raster embed \\
Cloud / atmosphere     & Not available        & Native                  & Raster embed \\
Per-pixel compositing  & Limited (opacity)    & Native                  & Limited \\
\midrule
\textit{Max fidelity tier} & Tier~1 (det.); Tier~2 (caveat) & Tier~3 & Tier~2 (native) $+$ Tier~3 (hybrid) \\
\bottomrule
\end{tabularx}
\caption{Backend capability profile by effect class; ``Non-det.\ filter'' means the SVG
  construct exists but is not deterministic across renderers/browsers}
\label{tab:backend-capability}
\end{table}

\subsection{Output Priority Recommendation}
\label{sec:output-priority}

The following priority ordering is recommended for implementation. The Scene architecture
means the first investment is in the Scene model itself, which unlocks all subsequent
backends.

\begin{enumerate}
  \item \textbf{Scene / Render IR model + SVG backend} --- \emph{Build first, always.}
        The Scene is the correctness foundation. The SVG backend is the fastest path to
        a visible, inspectable output: every layout defect is immediately apparent in SVG
        before any rasterisation or conversion step can obscure it. It also serves as the
        primary debugging surface during development. Targets Tier-0 and Tier-1 themes
        (Minimal, Consulting, Release) with full byte-determinism.

  \item \textbf{PNG via SVG backend} --- \emph{Immediate universal value.}
        Rasterise the SVG backend output at caller-specified DPI using a headless SVG
        rasteriser with embedded font support (\texttt{resvg} or \texttt{librsvg}).
        Unlocks the most common real-world use case: pasting a timeline into a document,
        email, or Slack message. The effort is minimal once SVG is correct.

  \item \textbf{PDF via SVG backend} --- \emph{Print and archival; consulting use case.}
        Convert the SVG backend output to PDF. The primary differentiator
        vs.\ Mermaid~\cite{mermaid2023}; essential for consulting deliverables.
        The SVG-to-PDF conversion path is well-understood and deterministic.
        Priority~3 because font subsetting and PDF metadata handling are needed but
        engineering effort is low relative to PPTX.

  \item \textbf{PPTX native-shape backend} --- \emph{Highest editability value.}
        Hybrid native-shapes-plus-raster-art-layer strategy. Approaches think-cell
        editability for Tier-0/1 vocabulary; exploits native PowerPoint effects for
        Tier-2. The most complex target; requires stable Scene geometry from earlier
        priorities.

  \item \textbf{Raster backend (Skia or Canvas)} --- \emph{Art-effect differentiator.}
        The enabling technology for Showcase and Keynote themes. Glow, bloom, cloud
        textures, volumetric atmospheres. Ship as an optional plugin or separate
        distribution target to keep the core distribution lean.

  \item \textbf{HTML / interactive} --- \emph{Developer experience and agent integration.}
        Nearly free once the SVG backend is correct (SVG-backed path) or the Raster
        backend is available (canvas-backed path). Essential for the VS~Code extension
        live preview and the MCP agent integration path described in
        Section~\ref{sec:agent-integration}.
\end{enumerate}

\subsection{Format Comparison Summary}
\label{sec:format-comparison}

\begin{table}[h]
\centering
\small
\begin{tabularx}{\textwidth}{lXXXXX}
\toprule
\textbf{Property} & \textbf{SVG} & \textbf{PNG} & \textbf{PDF} & \textbf{PPTX} & \textbf{HTML} \\
\midrule
Scalable               & Yes          & No           & Yes          & Yes           & Yes (SVG path) \\
Byte-deterministic     & Tier~0/1     & Per backend  & Per backend  & Geometry only & Geometry only \\
Art effects            & Tier~1 only  & Tier~3 (raster) & Both paths & Hybrid       & Both paths \\
Editable shapes        & Via Inkscape & No           & No           & Yes (native)  & No \\
Text searchable        & Yes          & No           & Yes          & Yes           & Yes \\
Print-ready            & Via viewer   & At 300~DPI   & Yes          & Via PPT       & No \\
Universal compat.      & High         & Very high    & Very high    & Office-only   & Browser-only \\
Interactive            & No           & No           & No           & No            & Yes \\
Impl.\ complexity      & Foundation   & Low          & Low--medium  & High          & Low \\
Build priority         & 1            & 2            & 3            & 4             & 6 \\
\bottomrule
\end{tabularx}
\caption{Output format properties and recommended build priority}
\label{tab:format-comparison}
\end{table}

\paragraph{Determinism across backends.}
The Scene geometry is byte-deterministic unconditionally: the layout pipeline contract
(Section~\ref{sec:rendering-determinism}) guarantees this regardless of which backend is
invoked. Each backend is deterministic given a pinned backend version and fixed effect
seeds; backend version is recorded in output metadata to make this reproducible. SVG and
PDF (via SVG backend) satisfy full byte-determinism for Tier-0/1. PNG satisfies
byte-determinism per pinned backend version. PPTX satisfies geometric determinism;
internal \texttt{r:id} values are derived from entity \texttt{id} fields rather than
auto-generated. HTML satisfies geometry determinism; its interactive JavaScript layer is
explicitly excluded from the determinism contract. Cross-backend pixel identity is
explicitly \emph{not} promised and is not a correctness requirement: it is expected and
correct for the SVG backend and the Raster backend to produce visually different outputs
from the same Scene when the theme requests effects beyond SVG's deterministic range.
This is a feature of the architecture, not a defect.
